\documentclass{article}
\usepackage{amsmath,amssymb,amsthm}
\usepackage{algorithm}
\usepackage[noend]{algpseudocode}
\usepackage{hyperref}
\usepackage{enumitem}
\usepackage{color}
\newcommand{\E}{\mathbb{E}}
\newcommand{\Var}{\mathrm{Var}}
\newcommand{\R}{\mathbb{R}}
\newcommand{\norm}[1]{\left\lVert #1\right\rVert}
\newtheorem{theorem}{Theorem}
\newtheorem{lemma}{Lemma}
\newtheorem{assumption}{Assumption}
\newtheorem{corollary}{Corollary}
\newtheorem{remark}{Remark}
\begin{document}

\section*{DP‑SGD with Gaussian Noise}

\section*{Mathematical Formulation}
Let $f:\R^{d}\to\R$ be a (possibly non‑convex) loss function defined as the average over $n$ examples,
\[
f(w)=\frac{1}{n}\sum_{i=1}^{n} f_i(w),\qquad w\in\R^{d}.
\]
At iteration $t\in\{0,1,\dots,T-1\}$ we draw a minibatch $\mathcal{B}_t\subset\{1,\dots,n\}$ of size $b$ uniformly at random,
compute the (clipped) stochastic gradient
\[
\tilde g_t=\frac{1}{b}\sum_{i\in\mathcal{B}_t}\operatorname{clip}\!\bigl(\nabla f_i(w_t),C\bigr),
\]
and add isotropic Gaussian noise
\[
g_t=\tilde g_t+ \xi_t,\qquad \xi_t\sim\mathcal N\!\bigl(0,\sigma^{2}I_d\bigr).
\]
The parameter update is
\[
w_{t+1}=w_t-\eta_t g_t,
\]
where $\eta_t>0$ is the learning rate.  
The noise scale $\sigma$ is chosen such that each iteration satisfies $(\varepsilon_0,\delta_0)$‑differential privacy; after $T$ iterations the overall privacy budget $(\varepsilon,\delta)$ is obtained via a composition accountant (e.g., moments accountant).

\section*{Pseudocode}
\begin{algorithm}[H]
\caption{DP‑SGD with Gaussian Noise}
\begin{algorithmic}[1]
\Require Learning rates $\{\eta_t\}_{t=0}^{T-1}$, clipping norm $C>0$, noise scale $\sigma>0$, minibatch size $b$, total iterations $T$.
\State Initialize $w_0\in\R^{d}$ (e.g. $w_0=0$).
\For{$t=0$ \textbf{to} $T-1$}
    \State Sample a minibatch $\mathcal{B}_t$ of size $b$ uniformly.
    \State Compute per‑example gradients $g_i^{(t)}=\nabla f_i(w_t)$ for $i\in\mathcal{B}_t$.
    \State \textbf{Clip:} $\displaystyle \hat g_i^{(t)}=
          \frac{g_i^{(t)}}{\max\{1,\|g_i^{(t)}\|/C\}}$.
    \State Aggregate: $\displaystyle \tilde g_t=\frac{1}{b}\sum_{i\in\mathcal{B}_t}\hat g_i^{(t)}$.
    \State Sample noise $\xi_t\sim\mathcal N(0,\sigma^{2}I_d)$.
    \State Form noisy gradient $g_t=\tilde g_t+\xi_t$.
    \State Update $w_{t+1}=w_t-\eta_t g_t$.
\EndFor
\State \textbf{Output:} $w_T$ (or the average $\bar w_T=\frac1T\sum_{t=1}^{T}w_t$).
\end{algorithmic}
\end{algorithm}

\section*{Dry Run Example}
Consider the one‑dimensional quadratic $f(w)=(w-3)^2$.
\[
\nabla f(w)=2(w-3).
\]
Set $w_0=0$, learning rate $\eta=0.1$, clipping norm $C=10$ (no clipping will occur), noise variance $\sigma^2=0$ (to illustrate the deterministic part), and batch size $b=1$ (single data point).  

\begin{center}
\begin{tabular}{c|c|c|c}
$t$ & $w_t$ & $g_t=\nabla f(w_t)$ & $w_{t+1}=w_t-\eta g_t$\\ \hline
0 & $0$ & $2(0-3)=-6$ & $0-0.1(-6)=0.6$\\
1 & $0.6$ & $2(0.6-3)=-4.8$ & $0.6-0.1(-4.8)=1.08$\\
2 & $1.08$ & $2(1.08-3)=-3.84$ & $1.08-0.1(-3.84)=1.464$\\
\end{tabular}
\end{center}
Corresponding objective values:
$f(0)=9$, $f(0.6)=5.76$, $f(1.08)=3.6864$, $f(1.464)=2.3490$.
The iterates move toward the optimum $w^\star=3$.

\newpage
\section*{Assumptions}
\begin{assumption}[Smoothness]\label{ass:smooth}
$f$ is $L$‑smooth, i.e., for all $x,y\in\R^{d}$,
\[
f(y)\le f(x)+\langle\nabla f(x),y-x\rangle+\frac{L}{2}\|y-x\|^{2}.
\]
\end{assumption}
\begin{assumption}[Bounded Gradient Norm after Clipping]\label{ass:clip}
For every $i$ and $w$, the clipped gradient satisfies $\|\operatorname{clip}(\nabla f_i(w),C)\|\le C$.
\end{assumption}
\begin{assumption}[Unbiased Minibatch Gradient]\label{ass:unbiased}
Conditioned on $w_t$, the averaged clipped gradient is an unbiased estimator of the true clipped gradient:
\[
\E_{\mathcal{B}_t}\bigl[\tilde g_t\mid w_t\bigr]=\frac{1}{n}\sum_{i=1}^{n}\operatorname{clip}\!\bigl(\nabla f_i(w_t),C\bigr)=:\bar g_t .
\]
\end{assumption}
\begin{assumption}[Noise Independence]\label{ass:noise}
The Gaussian noise $\xi_t$ is independent of the minibatch and of past randomness, with
\[
\E[\xi_t]=0,\qquad \E\bigl[\|\xi_t\|^{2}\bigr]=d\sigma^{2}.
\]
\end{assumption}
\begin{assumption}[Learning Rate]\label{ass:lr}
The stepsizes are constant $\eta_t\equiv\eta$ satisfying $0<\eta\le \frac{1}{L}$.
\end{assumption}

\section*{Main Convergence Theorem}
\begin{theorem}\label{thm:main}
Assume \ref{ass:smooth}–\ref{ass:lr} hold. Let $\{w_t\}_{t=0}^{T}$ be generated by DP‑SGD with Gaussian noise of variance $\sigma^{2}$. Then for the averaged iterate
\[
\bar w_T:=\frac{1}{T}\sum_{t=0}^{T-1}w_t,
\]
the expected squared gradient norm satisfies
\[
\frac{1}{T}\sum_{t=0}^{T-1}\E\!\bigl[\|\nabla f(w_t)\|^{2}\bigr]
\;\le\;
\underbrace{\frac{2\bigl(f(w_0)-f^\star\bigr)}{\eta T}}_{\text{optimization term}}
\;+\;
\underbrace{L\eta C^{2}}_{\text{clipping bias}}
\;+\;
\underbrace{L\eta d\sigma^{2}}_{\text{privacy noise term}} .
\tag{1}
\]
Consequently, choosing $\eta=\Theta\bigl(1/\sqrt{T}\bigr)$ yields
\[
\frac{1}{T}\sum_{t=0}^{T-1}\E\!\bigl[\|\nabla f(w_t)\|^{2}\bigr]
=O\!\Bigl(\frac{1}{\sqrt{T}}+L C^{2}\frac{1}{\sqrt{T}}+L d\sigma^{2}\frac{1}{\sqrt{T}}\Bigr).
\]
\end{theorem}

\section*{Theoretical Properties}
\begin{itemize}
\item \textbf{Stability.} The presence of Gaussian noise guarantees that the iterates remain bounded in expectation because the noise variance is finite and the step size is restricted by $1/L$.
\item \textbf{Hyper‑parameter Sensitivity.}
    \begin{enumerate}[label=\alph*)]
    \item \emph{Clipping norm $C$:} larger $C$ reduces the clipping bias term $L\eta C^{2}$ but may increase the privacy loss (requires larger $\sigma$).
    \item \emph{Noise scale $\sigma$:} appears linearly in the bound via $L\eta d\sigma^{2}$; smaller $\sigma$ improves utility but weakens privacy.
    \item \emph{Learning rate $\eta$:} appears both as $\frac{1}{\eta T}$ and multiplied by $C^{2},d\sigma^{2}$. The optimal $\eta$ balances these contributions, giving the $\Theta(1/\sqrt{T})$ scaling.
    \end{enumerate}
\item \textbf{Comparison to non‑private SGD.} Setting $\sigma=0$ and $C=\infty$ (no clipping) reduces (1) to the classic non‑convex SGD bound $\frac{2(f(w_0)-f^\star)}{\eta T}+L\eta G^{2}$ where $G$ bounds the stochastic gradient norm.
\end{itemize}

\section*{Discussion}
The theorem shows that DP‑SGD attains the same $O(1/\sqrt{T})$ decay of the average squared gradient norm as standard SGD, up to additive terms that reflect the privacy‑induced noise and the clipping operation. When the privacy budget is tight (large $\sigma$), the noise term dominates and the algorithm behaves like stochastic gradient descent with high variance. Conversely, with a generous privacy budget the bound approaches the non‑private rate.

\newpage
\section*{Preliminaries (Definitions and Lemmas)}
\begin{lemma}[Smoothness Inequality]\label{lem:smooth}
If $f$ is $L$‑smooth then for any $x,y\in\R^{d}$,
\[
f(y)\le f(x)+\langle\nabla f(x),y-x\rangle+\frac{L}{2}\|y-x\|^{2}.
\]
\end{lemma}
\begin{proof}
Standard consequence of $L$‑Lipschitz gradient; see e.g. \cite{Nesterov2004}.
\end{proof}

\begin{lemma}[Bounded Variance of Noisy Gradient]\label{lem:variance}
For the noisy gradient $g_t=\tilde g_t+\xi_t$,
\[
\E\!\bigl[\|g_t\|^{2}\mid w_t\bigr]\le \underbrace{\E\!\bigl[\|\tilde g_t\|^{2}\mid w_t\bigr]}_{\le C^{2}}+d\sigma^{2}.
\]
\end{lemma}
\begin{proof}
Using independence of $\xi_t$ and $\tilde g_t$,
\[
\E\!\bigl[\|g_t\|^{2}\mid w_t\bigr]
 =\E\!\bigl[\|\tilde g_t\|^{2}\mid w_t\bigr]
   +\E\!\bigl[\|\xi_t\|^{2}\bigr]
   +2\E\!\bigl[\langle\tilde g_t,\xi_t\rangle\mid w_t\bigr].
\]
The cross term vanishes because $\E[\xi_t]=0$ and $\xi_t$ is independent of $\tilde g_t$, yielding the claim. By clipping, $\|\tilde g_t\|\le C$, hence the bound.
\end{proof}

\section*{Main Proof (Step‑by‑Step Derivation)}
We prove Theorem~\ref{thm:main}. All expectations are taken over the randomness of minibatch sampling and Gaussian noise.

\begin{enumerate}
\item[Step 1.] Apply Lemma~\ref{lem:smooth} with $x=w_t$, $y=w_{t+1}=w_t-\eta g_t$:
\[
f(w_{t+1})\le f(w_t)-\eta\langle\nabla f(w_t),g_t\rangle+\frac{L\eta^{2}}{2}\|g_t\|^{2}.
\tag{2}
\]

\item[Step 2.] Take conditional expectation w.r.t. the randomness at iteration $t$ (minibatch and noise) while fixing $w_t$:
\[
\E\!\bigl[f(w_{t+1})\mid w_t\bigr]\le f(w_t)-\eta\,\E\!\bigl[\langle\nabla f(w_t),g_t\rangle\mid w_t\bigr]
+\frac{L\eta^{2}}{2}\E\!\bigl[\|g_t\|^{2}\mid w_t\bigr].
\tag{3}
\]

\item[Step 3.] Decompose $g_t=\tilde g_t+\xi_t$ and use linearity of expectation:
\[
\E\!\bigl[\langle\nabla f(w_t),g_t\rangle\mid w_t\bigr]
=\langle\nabla f(w_t),\E[\tilde g_t\mid w_t]\rangle
+\langle\nabla f(w_t),\E[\xi_t]\rangle.
\]
By Assumption~\ref{ass:noise}, $\E[\xi_t]=0$, and by Assumption~\ref{ass:unbiased},
\[
\E[\tilde g_t\mid w_t]=\bar g_t:=\frac{1}{n}\sum_{i=1}^{n}\operatorname{clip}(\nabla f_i(w_t),C).
\]
Hence
\[
\E\!\bigl[\langle\nabla f(w_t),g_t\rangle\mid w_t\bigr]
=\langle\nabla f(w_t),\bar g_t\rangle.
\tag{4}
\]

\item[Step 4.] Relate $\bar g_t$ to the true gradient. By definition of clipping,
\[
\operatorname{clip}(u,C)=u\quad\text{if }\|u\|\le C,\qquad
\operatorname{clip}(u,C)=\frac{C}{\|u\|}u\quad\text{otherwise}.
\]
Thus $\|\operatorname{clip}(u,C)-u\|\le \|u\|\,\mathbf 1_{\{\|u\|>C\}}$ and consequently
\[
\bigl\|\bar g_t-\nabla f(w_t)\bigr\|
\le\frac{1}{n}\sum_{i=1}^{n}\bigl\|\operatorname{clip}(\nabla f_i(w_t),C)-\nabla f_i(w_t)\bigr\|
\le\frac{1}{n}\sum_{i=1}^{n}\|\nabla f_i(w_t)\|\mathbf 1_{\{\|\nabla f_i(w_t)\|>C\}}.
\]
We do not need a finer bound; we will upper‑bound the inner product directly.

\item[Step 5.] Use the identity $2\langle a,b\rangle=\|a\|^{2}+\|b\|^{2}-\|a-b\|^{2}$ with $a=\nabla f(w_t)$ and $b=\bar g_t$:
\[
\langle\nabla f(w_t),\bar g_t\rangle
=\frac{1}{2}\Bigl(\|\nabla f(w_t)\|^{2}+\|\bar g_t\|^{2}-\|\nabla f(w_t)-\bar g_t\|^{2}\Bigr).
\tag{5}
\]

\item[Step 6.] Since $\|\bar g_t\|\le C$ by clipping, we have $\|\bar g_t\|^{2}\le C^{2}$. Moreover, the last term $-\|\nabla f(w_t)-\bar g_t\|^{2}\le 0$. Hence from (5)
\[
\langle\nabla f(w_t),\bar g_t\rangle
\ge\frac{1}{2}\bigl(\|\nabla f(w_t)\|^{2}-C^{2}\bigr).
\tag{6}
\]

\item[Step 7.] Plug (6) into (4) and then into (3):
\[
\E\!\bigl[f(w_{t+1})\mid w_t\bigr]
\le f(w_t)-\frac{\eta}{2}\|\nabla f(w_t)\|^{2}
+\frac{\eta}{2}C^{2}
+\frac{L\eta^{2}}{2}\E\!\bigl[\|g_t\|^{2}\mid w_t\bigr].
\tag{7}
\]

\item[Step 8.] Apply Lemma~\ref{lem:variance} to bound the last term:
\[
\E\!\bigl[\|g_t\|^{2}\mid w_t\bigr]\le C^{2}+d\sigma^{2}.
\tag{8}
\]

\item[Step 9.] Substitute (8) into (7):
\[
\E\!\bigl[f(w_{t+1})\mid w_t\bigr]
\le f(w_t)-\frac{\eta}{2}\|\nabla f(w_t)\|^{2}
+\frac{\eta}{2}C^{2}
+\frac{L\eta^{2}}{2}\bigl(C^{2}+d\sigma^{2}\bigr).
\tag{9}
\]

\item[Step 10.] Take total expectation over all randomness up to iteration $t$ and rearrange:
\[
\frac{\eta}{2}\E\!\bigl[\|\nabla f(w_t)\|^{2}\bigr]
\le \E\!\bigl[f(w_t)-f(w_{t+1})\bigr]
+\frac{\eta}{2}C^{2}
+\frac{L\eta^{2}}{2}\bigl(C^{2}+d\sigma^{2}\bigr).
\tag{10}
\]

\item[Step 11.] Sum (10) over $t=0,\dots,T-1$:
\[
\frac{\eta}{2}\sum_{t=0}^{T-1}\E\!\bigl[\|\nabla f(w_t)\|^{2}\bigr]
\le f(w_0)-\E\!\bigl[f(w_T)\bigr]
+\frac{T\eta}{2}C^{2}
+\frac{TL\eta^{2}}{2}\bigl(C^{2}+d\sigma^{2}\bigr).
\tag{11}
\]

\item[Step 12.] Since $f(w_T)\ge f^\star$, drop the negative term $-\E[f(w_T)]$:
\[
\frac{\eta}{2}\sum_{t=0}^{T-1}\E\!\bigl[\|\nabla f(w_t)\|^{2}\bigr]
\le f(w_0)-f^\star
+\frac{T\eta}{2}C^{2}
+\frac{TL\eta^{2}}{2}\bigl(C^{2}+d\sigma^{2}\bigr).
\tag{12}
\]

\item[Step 13.] Divide both sides by $T\eta$ and use $0<\eta\le 1/L$ so that $L\eta\le 1$:
\[
\frac{1}{T}\sum_{t=0}^{T-1}\E\!\bigl[\|\nabla f(w_t)\|^{2}\bigr]
\le \frac{2\bigl(f(w_0)-f^\star\bigr)}{\eta T}
+ C^{2}
+ L\eta\bigl(C^{2}+d\sigma^{2}\bigr).
\tag{13}
\]

\item[Step 14.] Rearranging the terms yields exactly the bound (1) stated in Theorem~\ref{thm:main}.
\end{enumerate}

\section*{Rate Derivation (With Constants)}
From (13) we identify three contributions:
\begin{itemize}
\item Optimization term: $\displaystyle\frac{2\bigl(f(w_0)-f^\star\bigr)}{\eta T}$.
\item Clipping bias term: $L\eta C^{2}$.
\item Privacy‑noise term: $L\eta d\sigma^{2}$.
\end{itemize}
Choosing $\eta = \sqrt{\frac{2\bigl(f(w_0)-f^\star\bigr)}{(C^{2}+d\sigma^{2})LT}}$ (which respects $\eta\le 1/L$ for sufficiently large $T$) balances the optimization term with the sum of the two linear terms, giving
\[
\frac{1}{T}\sum_{t=0}^{T-1}\E\!\bigl[\|\nabla f(w_t)\|^{2}\bigr]
= O\!\Bigl(\sqrt{\frac{L\bigl(C^{2}+d\sigma^{2}\bigr)}{T}}\Bigr).
\]
In the regime where $C$ and $d\sigma^{2}$ are constants, this simplifies to the familiar $O(1/\sqrt{T})$ rate.

\section*{Special Cases}
\begin{enumerate}
\item \textbf{No privacy noise ($\sigma=0$).} The bound reduces to the standard non‑convex SGD result with clipping:
\[
\frac{1}{T}\sum_{t=0}^{T-1}\E\!\bigl[\|\nabla f(w_t)\|^{2}\bigr]
\le \frac{2\Delta_f}{\eta T}+L\eta C^{2},
\quad \Delta_f:=f(w_0)-f^\star .
\]
\item \textbf{No clipping ($C\to\infty$) but with noise.} Then $C^{2}$ terms vanish, and we obtain
\[
\frac{1}{T}\sum_{t=0}^{T-1}\E\!\bigl[\|\nabla f(w_t)\|^{2}\bigr]
\le \frac{2\Delta_f}{\eta T}+L\eta d\sigma^{2},
\]
which matches the analysis of pure Gaussian‑perturbed SGD.
\item \textbf{Convex smooth case.} If $f$ is convex, the same derivation yields a bound on suboptimality:
\[
\E\!\bigl[f(\bar w_T)-f^\star\bigr]\le
\frac{\|w_0-w^\star\|^{2}}{2\eta T}
+ \frac{L\eta}{2}\bigl(C^{2}+d\sigma^{2}\bigr).
\]
\end{enumerate}

\begin{thebibliography}{9}
\bibitem{Nesterov2004}
Y.~Nesterov, \emph{Introductory Lectures on Convex Optimization: A Basic Course}, Kluwer Academic Publishers, 2004.
\end{thebibliography}

\end{document}